\section{DA as estimation of a conditional distribution}\label{sec:chain}

\subsection{State-space model and context}
Let $\xv = (\xv_0, \dots, \xv_K)$ be the state trajectory over an assimilation
window, $\zv$ an external forcing and $\thv$ the physical parameters. The
discrete dynamics are
\begin{equation}\label{eq:dynamics}
  \xv_{k+1} = \Mop_\Delta(\xv_k, \zv_k; \thv) + \xi_k ,
\end{equation}
with $\Mop_\Delta$ the time-stepping operator over a physical step $\Delta$ and
$\xi_k$ a model-error process; the observations are
\begin{equation}\label{eq:obs}
  \yv_k = \Hop(\xv_k) + \epsv_k, \qquad \epsv_k \sim \mathcal{N}(0, \mathbf{R}),
\end{equation}
at the observation times of the window. Everything an estimator may condition on
is collected in the \emph{context} $C = (\yv, \thv, \zv)$. The posterior
\begin{equation}\label{eq:posterior}
  p(\xv \mid C) \;\propto\; p(\yv \mid \xv)\; p(\xv \mid \thv, \zv)
\end{equation}
combines the observation likelihood (through $\Hop$ and $\mathbf{R}$) with a prior
that, in model-based DA, is generated by the dynamics~\eqref{eq:dynamics}.
Learned schemes replace or bypass this prior. \textbf{The target throughout is the
state alone}: $\thv$ and $\zv$ are inputs that a scheme may use, never
quantities it estimates.

The conditional expectation $m = \E[\xv \mid C]$ is the MSE-optimal estimator,
and $\E \tr \Cov(\xv \mid C)$ is the smallest mean-squared error any scheme can
reach: the \emph{irreducible} error.

\subsection{The error chain}
A DA scheme departs from $m$ in three ways. It uses some information
$S \subseteq C$ --- a filter uses the observations up to the current time, a
learned scheme may ignore $\thv$ and $\zv$. It believes a model $\tilde p$ that may
differ from the true one. And it computes its own target only approximately. We
write the corresponding chain of estimators as
\begin{equation}\label{eq:chain}
  m = \E[\xv\mid C] \;\longrightarrow\; m_S = \E[\xv\mid S]
  \;\longrightarrow\; \tilde m = \E_{\tilde p}[\xv\mid S]
  \;\longrightarrow\; \hat m ,
\end{equation}
where $\hat m$ is what the scheme outputs. If $\hat m$ depends only on $S$ and on
randomness independent of the truth (ensemble draws, initial weights), then
\begin{equation}\label{eq:chain0}
  \mathrm{MSE}(\hat m) \;=\;
  \underbrace{\E\tr\Cov(\xv\mid C)}_{\text{irreducible}}
  \;+\; \underbrace{\E\|m - m_S\|^2}_{\text{restriction}}
  \;+\; \underbrace{\E\|m_S - \hat m\|^2}_{\text{misspecification + approximation}} .
\end{equation}
Both splits are exact. The first holds because $m$ is the conditional expectation
given $C$. The second holds because $m_S = \E[m \mid S]$ (tower property), while
$m_S - \hat m$ is a function of $S$ and of independent randomness only, so the
cross terms vanish.

The last term contains the misspecification $\E\|m_S - \tilde m\|^2$, the
approximation $\E\|\tilde m - \hat m\|^2$ and their cross term. That cross term
does not vanish in general: a finite ensemble can partly compensate a wrong
model, or worsen it. We therefore do not claim a four-term identity. We separate
misspecification from approximation by controlled contrasts --- the same scheme
with more members, more iterations or more training, and the same scheme with
and without model error (Sect.~\ref{sec:design}) --- and we report the cross term
as the residual it is.

\subsection{Where the commitments of DA land}
Each commitment a DA scheme makes lands on one link of the chain
(Table~\ref{tab:commitments}).
\begin{itemize}[leftmargin=*]
  \item \textbf{Restriction.} Filtering conditions $\xv_k$ on $\yv_{t \le k}$ only. Cycling a variational
    scheme over sub-windows carries later observations only through a static
    background. Learned schemes that ignore $\thv$ and $\zv$ marginalise over
    their training spread.
  \item \textbf{Misspecification.} Model error in the dynamics, parameters and forcing (our \So{} case)
    affects every scheme that consumes the model. Gaussian updates, the mode
    instead of the mean, a hard constraint ($\mathbf{Q} = 0$) and the
    tempered likelihood of guided samplers are structural misspecifications that
    remain even with a perfect model.
  \item \textbf{Approximation.} Finite ensembles, local optimisation of
    non-convex costs, and finite networks and training data make the computed
    output differ from the scheme's own target.
\end{itemize}

\begin{table}[t]
\centering\small
\caption{Where the commitments of DA and of learned schemes land in the error chain~\eqref{eq:chain}.}
\label{tab:commitments}
\begin{tabular}{L{5.6cm}L{2.7cm}L{5.6cm}}
\toprule
commitment & link & schemes concerned \\
\midrule
filtering ($S = \yv_{t'\le t}$); cycled sub-windows & restriction & EnKF, ETKF (vs ETKS); 4D-Var \\
ignoring $\thv$, $\zv$ & restriction & DirectUNet, CFM, SDA1 \\
model error in $\Mop$, $\thv$, $\zv$ & misspecification & every model-based scheme at \So \\
Gaussian update; mode instead of mean; hard constraint; tempered likelihood & misspecification (structural) & ETKF/EnKF; Strong-4DVar; SDA guidance \\
finite ensemble; local optimisation; finite network and data & approximation & $N = 30$; L-BFGS; 1200 epochs \\
\bottomrule
\end{tabular}
\end{table}

These commitments are choices, not errors: each buys something --- tractability,
a sequential implementation, robustness to the observing system --- and the
question is what it costs, and when the cost changes nature. The five design
choices behind them (prior--likelihood split, source of the dynamics, sequential
inference, how state uncertainty and how model uncertainty are represented)
label every scheme in Table~\ref{tab:schemes} and are defined in
Appendix~\ref{app:matrix}.
